\section{AlphaCode: búsqueda a gran escala en el espacio de programas}

\subsection{Por qué la programación competitiva es más difícil que el completado de código}

HumanEval suele dar la firma de la función, una descripción local y una implementación corta; los problemas de Codeforces, en cambio, exigen entender un enunciado extenso, elegir un algoritmo, satisfacer restricciones de complejidad y escribir un programa completo. El error puede surgir en cualquier eslabón: comprensión del enunciado, algoritmo, casos límite, complejidad, sintaxis o el protocolo de entrada y salida.

\videofigure{figures/alphacode-pipeline.jpg}{AlphaCode encadena en un solo flujo el preentrenamiento de código, el fine-tuning con problemas de competencia, el muestreo a gran escala, el filtrado y la agrupación, y el envío final.}{00:03:48--00:05:42}

Los datos de entrenamiento provienen principalmente de código de GitHub y de CodeContests. El AlphaCode original todavía requiere construir su propio modelo: el preentrenamiento aprende la distribución general del código, y el fine-tuning con datos de competencia adapta el modelo a la estructura «enunciado--solución»; solo después entra en la búsqueda en tiempo de prueba.

\begin{warningbox}{El loss del conjunto de validación no es un indicador confiable del solve rate}
El loss a nivel de token mide qué tan cerca está la predicción local de la distribución de entrenamiento, pero los problemas de competencia se califican según si el programa completo pasa las pruebas ocultas. Un error de límite en un solo token puede anular todo el problema; por eso una mejora pequeña en el loss de entrenamiento no necesariamente se traduce en una mejora en la tasa de problemas resueltos.
\end{warningbox}

\subsection{Muestrear un millón de programas por problema}

\videofigure{figures/large-scale-sampling.jpg}{Para cada problema, AlphaCode genera cerca de un millón de candidatos en Python/C++, y usa temperatura alta y perturbaciones del prompt para mantener la diversidad algorítmica.}{00:06:08--00:07:28}

Sea $p_\theta(y\mid x)$ la distribución condicional del modelo base sobre el programa $y$. Si la probabilidad de que un solo muestreo acierte con el programa correcto es $p$, la probabilidad de acertar al menos una vez en $n$ muestreos independientes es:

$$
P(\text{hit})=1-(1-p)^n.
$$

\begin{itemize}
    \item $x$: el enunciado, los ejemplos y las etiquetas opcionales;
    \item $y$: el programa candidato completo;
    \item $p$: la probabilidad de que un solo muestreo caiga en el conjunto de soluciones correctas;
    \item $n$: el presupuesto de muestreo.
\end{itemize}

Cuando $p$ es extremadamente pequeña, el muestreo a gran escala puede aumentar significativamente la cobertura; pero la ganancia termina limitada por el selector, porque la plataforma solo permite enviar muy pocos programas.

\begin{importantbox}{El objetivo de la temperatura alta es cubrir familias de algoritmos, no generar diferencias sintácticas al azar}
La diversidad ideal proviene de enfoques de solución distintos: programación dinámica, algoritmos voraces, algoritmos de grafos, estructuras de datos, entre otros. Si los candidatos solo difieren en los nombres de variables o son reescrituras de la misma plantilla errónea, un número mayor de muestras no aumenta la cobertura efectiva.
\end{importantbox}

\subsection{Primero pasar los ejemplos, luego agrupar por semántica}

\videofigure{figures/filter-cluster.jpg}{El sistema primero filtra los programas que no pasan los ejemplos públicos, y luego agrupa en un mismo clúster los programas semánticamente similares según su comportamiento sobre entradas de prueba generadas.}{00:07:28--00:09:04}

AlphaCode entrena además un test-input generator. Para el mismo problema, genera un conjunto de entradas $T=\{t_j\}$ y mapea el programa $y$ a una firma de comportamiento:

$$
\phi_T(y)=\bigl(y(t_1),y(t_2),\ldots,y(t_m)\bigr).
$$

Si las firmas de comportamiento de dos programas son parecidas, es posible que implementen el mismo algoritmo o que compartan el mismo error; tras la agrupación, conservar solo unos cuantos representantes de cada clúster permite cubrir más semánticas distintas dentro de un presupuesto de envíos fijo.

\begin{knowledgebox}{La agrupación es un asignador de presupuesto}
El número de envíos en la plataforma es un recurso escaso. La agrupación no busca que los programas se vean ordenados, sino evitar que los diez envíos se desperdicien en el mismo patrón de fallo.
\end{knowledgebox}

\subsection{Resumen de la sección}

La clave de AlphaCode es comprimir un espacio de programas enorme en un número reducido de candidatos diversos: el modelo se encarga de la propuesta, los ejemplos públicos hacen un filtrado barato, las pruebas generadas realizan la agrupación por comportamiento, y las pruebas ocultas completan la verificación final.
